\section{Black-box Evaluations and Stopping Criteria}\label{sec:conv}

\subsection{Black-box Evaluations}\label{subsec:eval}

In the context of the frameworks presented in this paper, one black-box evaluation refers to a single complete execution of a disciplinary simulation to map a specific set of inputs to the output vector \(\bm{y}_i\). It is designated as a black-box because the optimization algorithm only observes this input-output numerical relationship and cannot exploit the underlying algebraic structure or exact internal analytical derivatives of the governing physical equations. This is typically the most computationally expensive operation in the overall design process.

In CO, the computational cost is fundamentally dictated by the evaluation frequency of \(\bm{y}_i\). Because the discrepancy formulation mandates that \(\bm{y}_i\) is evaluated as a explicit function of the active system-level coupling variables, \(\ysys_{j\neq i}\), these expensive function calls cannot be isolated strictly to the lower-level subsystem optimizations. Every iteration of the upper system-level optimizer that perturbs the target states \(\ysys\) forces a simultaneous re-evaluation of \(\bm{y}_i\) across all \(N\) subsystems to accurately compute the discrepancy constraint \(J\). Letting \(K\) denote the total number of major iterations to reach convergence, \(L\) the average number of system-level optimization iterations per major iteration, and \(M_i\) the average number of local evaluations performed by the \(i\)-th subsystem optimizer per major iteration, the total computational complexity in terms of black-box calls scales as \(\Ord\left(K \left( N L + \sum_{i=1}^{N} M_i \right)\right)\). This strict coupling indicates that for expensive-to-compute multidisciplinary problems, the computational burden is heavily amplified by the dimensionality of the coupling variables, as the upper-level search directly multiplies the required number of high-fidelity simulation runs.

In BACO, nested optimization searches are replaced with GP surrogate evaluations to reduce computational cost. During each major iteration \(k\), the subsystem-level optimizers maximize an acquisition function \(\alpha_{J_i}\) subject to surrogate constraint predictions \(\bm{\mu}_{g_i}\). This internal search relies exclusively on surrogates, requiring no evaluations of the true disciplinary black-box functions \(\bm{y}_i\). Instead, \(\bm{y}_i\) is evaluated exactly once per subsystem after local optimization terminates to establish the true discrepancy and constraint values, which update the data sets \(\underline{\D_{J_i}}^{(k+1)}\) and \(\D_i^{(k+1)}\). Similarly, the system-level optimizer manipulates the coupling variables \(\ysys\) using the surrogate models \(\alpha_f\), \(\bm{\mu}_c\), and \(\mu_{J_i}\). Following this system-level update, \(\bm{y}_i\) is evaluated a second time across all \(N\) subsystems to determine the true discrepancy at the new system state, augmenting the upper-level data set \(\overline{\D_{J_i}}^{(k+1)}\). Because all iterative search procedures are performed on the GP surrogates, the framework strictly dictates two disciplinary evaluations per discipline per major iteration. For \(K\) major iterations, the exact total number of true black-box calls is \(2KN\), corresponding to an asymptotic computational complexity of \(\Ord(KN)\).

\subsection{Discrepancy and Constraint Violation Metrics}\label{subsec:metrics}

The state of a CO iteration is characterized by two scalar metrics evaluated at the system level after each subsystem optimization cycle. For a given system-level iterate~\((\zsys, \xsys, \ysys, \zsub^*, \xsub^*)\), the total discrepancy \(\Jtot\) aggregates the inter-disciplinary incompatibility across all subsystems,
\begin{equation}
  \Jtot\left(
    \zsys,
    \xsys,
    \ysys,
    \zsub^*,
    \xsub^*
  \right) = \sum_{i=1}^N J\left(
    \zsys,
    \xsys_i,
    \ysys,
    \zsub_i^*,
    \xsub_i^*
  \right),\label{eqn:disc_total}
\end{equation}
while the total constraint violation~\cite{audet2017}, \(\htot\), aggregates both subsystem-level and system-level infeasibility,
\begin{equation}
  \htot\left(
    \zsys,
    \xsys,
    \ysys,
    \zsub^*,
    \xsub^*
  \right)
  = \sum_{i=1}^N \sum_{l=1}^{m_i} \max\left\{
    0, - g_{i_l}\left(
      \ysys_{j\neq i},
      \zsub_i^*,
      \xsub_i^*
    \right)
  \right\}^{2}
  + \sum_{l=1}^m \max\left\{
    0, - c_l\left(
      \zsys,
      \xsys,
      \ysys
    \right)
  \right\}^{2}.\label{eqn:constviol_total}
\end{equation}
Both metrics are non-negative by construction, with \(\Jtot = 0\) attained only at a fully compatible point and \(\htot = 0\) only at a feasible point.

These metrics serve a dual role. They provide a stopping criterion. The iteration terminates at step \(k\) when \(\Jtot(\zsys^{(k)}, \ldots) \leq \epsilon_J\) and \(\htot(\zsys^{(k)}, \ldots) \leq \epsilon_h\) for prescribed tolerances \(\epsilon_J, \epsilon_h > 0\). They also serve as the primary axes along which solution quality is assessed, complementing the objective value \(f\).
